\chapter{DPU-Accelerated Post-Quantum Cryptography}
\label{ch:dpu-pqc}

\section{Post-Quantum Cryptography and Its Deployment Cost}
\label{sec:pqc-background}

Post-quantum cryptography (PQC) represents the computationally secure counterpart to QKD.
It replaces the integer factorization and discrete logarithm problems vulnerable to Shor's
algorithm with entirely different mathematical structures believed to resist both classical
and quantum cryptanalysis. The transition focuses on two main primitive types: key
encapsulation mechanisms (KEMs), which establish a shared symmetric key over an insecure
channel, and digital signatures, which provide authentication and non-repudiation.

\subsection{Algorithmic Families and Characteristics}

The standardized and candidate algorithms studied in this part vary drastically in their
underlying mathematics, leading to profound differences in computational structure,
memory footprint, and message sizes~\cite{nistfips203,nistfips204,nistfips205,nistir8545}.

\begin{itemize}
    \item \textbf{Lattice-Based Cryptography:} Algorithms like ML-KEM (formerly Kyber) and
    ML-DSA (formerly Dilithium) rely on the hardness of the module learning with errors (MLWE)
    problem. They offer an excellent balance: fast execution times, moderate key and ciphertext
    sizes, and efficient hardware implementations. However, their reliance on polynomial sampling
    and rejection sampling introduces variable timing that must be carefully managed to avoid
    side-channel leaks.
    \item \textbf{Code-Based Cryptography:} Algorithms like HQC (Hamming Quasi-Cyclic) rely on
    the hardness of decoding random linear codes. While their encryption and decryption operations
    are fast, they suffer from significantly larger public keys and ciphertexts compared to lattices.
    Additionally, their operations involve large, dense bit-vector multiplications that have distinct
    memory access patterns~\cite{meyer2026hqc}.
\end{itemize}

Table~\ref{tab:pqc-comparison} summarizes these structural differences.

\begin{table}[htbp]
\centering
\caption{Comparison of selected Post-Quantum Cryptographic primitives and their deployment profiles. Sizes and operations are approximations based on NIST security category 3~\cite{nistfips203,nistfips204,nistir8545}.}
\label{tab:pqc-comparison}
\renewcommand{\arraystretch}{1.3}
\footnotesize
\begin{tabularx}{\textwidth}{llcrr>{\raggedright\arraybackslash}X}
\toprule
\textbf{Algorithm} & \textbf{Family} & \textbf{Type} & \textbf{PK (B)} & \textbf{CT/Sig (B)} & \textbf{Compute profile} \\
\midrule
ML-KEM-768 & Lattice & KEM & 1,184 & 1,088 & Fast polynomial arithmetic; SIMD-friendly (AVX/Neon). \\
ML-DSA-65 & Lattice & Signature & 1,952 & 3,309 & Rejection sampling; highly variable timing. \\
HQC-192 & Code & KEM & 4,522 & 9,026 & Bitwise math; memory-bound (large keys). \\
\bottomrule
\end{tabularx}
\end{table}

\subsection{Deployment and Architectural Constraints}

For networked systems, the migration to PQC is not finished when a NIST standard is published.
A mechanism that executes efficiently in a single-threaded benchmark may behave disastrously
under datacenter provisioning constraints.

Integrating PQC into protocols like TLS 1.3 or IPsec requires transmitting much larger
public keys and ciphertexts during handshakes. This can trigger IP fragmentation, increasing
latency and packet-loss probability. Furthermore, datacenters handle thousands of connections
per second. KEM decapsulation and signature verification must therefore be deeply batched.
Lattice schemes, with their variable-time rejection sampling, batch poorly on SIMD architectures
because a single long-running sample stalls the entire vector lane.

Consequently, the central questions are architectural: where should the cryptographic work run,
how should requests be batched, how much latency can the service mesh tolerate, and which
memory buses become bottlenecks? These are the questions that motivate offloading PQC onto
the DPU substrate introduced in Chapter~\ref{ch:background}, and they organize the two
contributions that follow.

\section{Contribution and Articles}

This chapter is built around two contributions that study how post-quantum cryptographic workloads can be executed efficiently on heterogeneous, programmable network hardware.

The first contribution (P4), \enquote{Offloading Post-Quantum Cryptography to Data Processing Units: A Performance Analysis of HQC Key Encapsulation}, analyzes the offload of HQC key encapsulation to a BlueField DPU and characterizes the resulting performance.
% TODO: Record the author's specific contributions and co-author context for P4.

The second contribution (P5), \enquote{GPU-Accelerated Post-Quantum Cryptography on Embedded DPU Architecture}, develops and evaluates a heterogeneous DPU architecture that places GPU compute directly in the network datapath to accelerate post-quantum cryptography across CPU, GPU, and DPU resources.
% TODO: Record the author's specific contributions and co-author context for P5.

Beyond these articles, the work on post-quantum cryptographic offload also contributed to a filed patent application, \enquote{Computational Offloading of Operations Necessary for Post-Quantum Cryptography on Network Interface Cards}, reported at title level among the patents in the List of Publications (page~\pageref{listPublications}).

TODO: Synthesize the findings across HQC offload and heterogeneous DPU acceleration, including what they reveal about bottlenecks, host interaction, and deployment constraints, and transition from isolated cryptographic acceleration to full application pipelines that combine data movement, AI workloads, and PQC.

\includedpaper{P4}{Offloading Post-Quantum Cryptography to Data Processing Units: A Performance Analysis of HQC Key Encapsulation}{meyer2026hqc}{source-materials/papers/camera-ready/hqc-pqc-offload-bluefield-dpu-ieee-access-2026-published.pdf}

\includedpaperpending{P5}{GPU-Accelerated Post-Quantum Cryptography on Embedded DPU Architecture}{meyer2026heterogeneous}
